\begin{afterword}
\summaryhead

The post extends ``Correspondence Bias'' to enemies. When someone offends us, the author observes, we are even more inclined to explain the act by an evil disposition. Prior probability says we should instead ask what the Enemy believes about the situation. Most people see themselves as heroes, so a portrait that makes the Enemy look bad will be wrong. Politics makes this hard: arguments are soldiers, and correcting a fact about the Enemy looks like treason.

Understanding is not approval. If an army invades or a lunatic attacks, killing someone who might have been your friend is a tragedy, but dying is also a tragedy, and wanting a non-tragic option is wanting a one-sided policy debate. An accurate picture of the Enemy should make you sad, not righteous. So the 9/11 hijackers were not evil mutants who hated freedom. They were heroes of their own stories, who in another environment might have been police officers.

\responsehead

The central claim is plausible and has support the post does not cite. A study of ordinary conflicts found that people who cause harm describe it as meaningful, while their victims describe it as gratuitous. On the hijackers, Osama bin Laden rejected the ``hate freedom'' account in 2004 and named political grievances instead. The post's distinction between understanding and approval is also kept with care: three sentences in the last paragraph say what the claim does not imply.

The argument has three weak points. First, the premise that the bias grows stronger when we are offended is supported only by the author's observation. Second, the view the post rejects is stated in a form no one is quoted holding. The ``hate freedom'' account, as President Bush gave it, explained the hatred by an extremist ideology, not by how the hijackers were born. Third, like ``Belief as Attire,'' the post asserts what the hijackers believed without citing anything they or their leaders said.

One passage does what the post warns against. Having warned that construing motives which make the Enemy look bad will leave you ``flat wrong,'' the post names Scientologists and ``Jesus Campers,'' without evidence, as people whose beliefs have ``utterly destroyed their sanity.'' It blames their beliefs, not their birth, so it does not call them innately evil. But it gives two whole groups the unflattering portrait the post warns against, without evidence.

In short: a sound warning against explaining enemies by innate evil, supported by evidence it does not cite, aimed at a version of the rival view no one is quoted holding, and not extended to two groups the post itself names.
\end{afterword}
